\section{Giới thiệu}
\label{sec:intro}

% Viết theo đoạn văn, hạn chế bullet (xem long-form best practices trong skill).
UAV làm relay là một hướng được quan tâm trong truyền thông 6G nhờ khả năng triển khai nhanh và thiết lập liên kết tầm nhìn thẳng với thiết bị mặt đất~\cite{huang2025ddatsap}. Trong bối cảnh cứu hộ sau thiên tai, yêu cầu quan trọng không phải là tốc độ trung bình mà là việc mỗi cảnh báo được giao đủ dữ liệu trước một hạn cho trước. [Đoạn mở rộng động lực nghiên cứu.]

Bài Huang và cộng sự~\cite{huang2025ddatsap} tối ưu đồng thời hệ số lập lịch tài nguyên, công suất và quỹ đạo để tối đa tốc độ trung bình nhỏ nhất. Đề tài này kế thừa mô hình kênh và cấu trúc phân rã của công trình đó nhưng thay đổi mục tiêu sang tỷ lệ cảnh báo đúng hạn và bổ sung lựa chọn relay. [Đoạn nêu khoảng trống.]

\paragraph{Đóng góp.}
Báo cáo có ba đóng góp chính. Thứ nhất, chúng tôi mô hình hoá bài toán với release time, deadline và mạng relay mặt đất có dung lượng hữu hạn (\cref{sec:method}). Thứ hai, chúng tôi đề xuất thuật toán phân rã theo khối kèm cơ chế sửa lịch theo slack, và đưa ra cận trên bằng nới lỏng tuyến tính để đo khoảng cách tối ưu. Thứ ba, chúng tôi đánh giá trên kịch bản tổng hợp từ dữ liệu công khai với 20--30 seed và kiểm định thống kê (\cref{sec:experiments}).

Phần còn lại của báo cáo được tổ chức như sau. \Cref{sec:related} điểm qua nghiên cứu liên quan. \Cref{sec:method} trình bày mô hình và thuật toán. \Cref{sec:experiments} báo cáo kết quả thực nghiệm, \cref{sec:discussion} thảo luận, và \cref{sec:conclusion} kết luận.
